% Auto-generated by review/make_supplementary.py
\documentclass[11pt]{article}
\usepackage[margin=1in]{geometry}
\usepackage{booktabs}
\usepackage{longtable}
\usepackage{amsmath}
\usepackage{times}
\usepackage{graphicx}
\graphicspath{
  {figures/}
  {figures/aesthetic/}
  {figures/aesthetic/families/}
}
\begin{document}
\begin{center}
{\Large\bfseries Supplementary Material}\\[6pt]
{\large Pitch--Hue Correspondences: A Quantitative Corpus Analysis of
Three Centuries of Mapping Systems}
\end{center}
\section*{S1. The full 70-system catalogue}
\begin{footnotesize}
\begin{longtable}{@{}p{5.6cm}p{1.6cm}p{2.6cm}p{1.4cm}p{4.6cm}@{}}
\caption{Table S1. The full 70-system catalogue: name, date, primary reasoning family, retention status, and retention note (mapping type; pitch coverage) or exclusion reason.}\label{tab:s1}\\
\toprule
\textbf{System} & \textbf{Date} & \textbf{Family} & \textbf{Status} & \textbf{Note} \\
\midrule\endfirsthead
\toprule
\textbf{System} & \textbf{Date} & \textbf{Family} & \textbf{Status} & \textbf{Note} \\
\midrule\endhead
\bottomrule\endfoot
Ancient Persian System & \textasciitilde{}700 BCE & F3 Cosmological & retained & pitch-class-to-color; diatonic; 7 keys \\
Chinese Wuxing Five-Tone/Five-Color System & \textasciitilde{}4th century BCE onward & F3 Cosmological & retained & pitch-class-to-color; pentatonic; 5 keys \\
Aristotle / Cureau de la Chambre Interval System & \textasciitilde{}350 BCE / 1650 CE & F3 Cosmological & discarded & interval-to-color; not pitch-class \\
Indian Navarasa Color System & \textasciitilde{}2nd century BCE & F3 Cosmological & discarded & rasa-to-color; not pitch-class \\
Indian Svara-to-Color Associations & Traditional, codified \textasciitilde{}12th century & F3 Cosmological & retained & pitch-class-to-color; heptatonic (saptak); 1 keys \\
Al-Kindi's Cosmological Oud-String System & \textasciitilde{}850 CE & F3 Cosmological & discarded & string/timbre-to-color; not pitch-class \\
Rudolph of Saint Trond Modal Colors & \textasciitilde{}1075 & F3 Cosmological & discarded & mode-to-color; not pitch-class \\
Franchino Gaffurio Modal Colors & 1492 & F3 Cosmological & discarded & mode-to-color; not pitch-class \\
Gioseffo Zarlino Interval-Color System & 1558 & F3 Cosmological & discarded & interval-to-color; not pitch-class \\
Girolamo Cardanus Interval-Color System & 1570 & F3 Cosmological & discarded & interval-to-color; not pitch-class \\
Athanasius Kircher Interval-Color System & 1646/1650 & F3 Cosmological & discarded & interval-to-color; not pitch-class \\
David Gottlob Diez Note-Color-Planet System & 1723 & F3 Cosmological & retained & pitch-class-to-color; chromatic (partial); 7 keys \\
Lorenz Christoph Mizler Newtonian System & 1739 & F1 Spectral Analogy & retained & pitch-class-to-color; diatonic; 7 keys \\
Johann Gottlob Krüger System & 1743 & F1 Spectral Analogy & retained & pitch-class-to-color; diatonic; 7 keys \\
Leonhard Euler System & \textasciitilde{}1760 & F1 Spectral Analogy & retained & pitch-class-to-color; diatonic; 7 keys \\
Louis François Henri Lefébure System & 1789 & F1 Spectral Analogy & retained & pitch-class-to-color; diatonic; 7 keys \\
Georg Tobias Ludwig Sachs (First Documented Synesthete) & 1812 & F5 Synesthetic & retained & pitch-class-to-color; partial; 5 keys \\
D.D. Jameson Chromatic Spectral Mapping & 1844 & F1 Spectral Analogy & retained & pitch-class-to-color; chromatic; 12 keys \\
Charles Fourier Inverted Spectral System & 1846 & F1 Spectral Analogy & retained & pitch-class-to-color; diatonic; 7 keys \\
François Sudre / Solresol System & 1862 & F8 Pedagogical & retained & pitch-class-to-color; diatonic; 7 keys \\
T. Seemann 12-Note System & 1881 & F1 Spectral Analogy & retained & pitch-class-to-color; chromatic; 12 keys \\
Bleuler and Lehmann Synesthesia Survey & 1881 & F5 Synesthetic & discarded & pitch-to-brightness/lightness only; not hue \\
Bainbridge Bishop Color Organ Spectral System & 1893 & F1 Spectral Analogy & retained & pitch-class-to-color; chromatic; 12 keys \\
Théodore Flournoy Synesthesia Cases & 1893 & F5 Synesthetic & discarded & qualitative case descriptions \\
Helena Blavatsky / Theosophical System & \textasciitilde{}1900 & F6 Esoteric & retained & pitch-class-to-color; diatonic; 7 keys \\
Rosicrucian Order System & \textasciitilde{}1905 & F6 Esoteric & retained & pitch-class-to-color; chromatic; 12 keys \\
Amy Beach Synesthetic Key-Colors & \textasciitilde{}1910 & F5 Synesthetic & retained & pitch-class-to-color; partial (key-based); 8 keys \\
Leonid Sabaneyev Documentation of Scriabin & 1911 & F2 Circle of Fifths & discarded & pitch-class-to-color; keys not normalizable to 12 pitch classes \\
Wassily Kandinsky Timbre-Color System & 1911 & F4 Emotion-Mediated & discarded & string/timbre-to-color; not pitch-class \\
Roy de Maistre Spectral System & 1919 & F1 Spectral Analogy & retained & pitch-class-to-color; chromatic; 12 keys \\
Leon Theremin / Illumovox Continuous System & 1922 & F1 Spectral Analogy & discarded & continuous/microtone; not discrete pitch-class \\
Adrian Bernard Klein Spectroscope System & 1926 & F1 Spectral Analogy & retained & pitch-class-to-color; chromatic; 12 keys \\
Artists Who Explicitly Rejected Fixed Pitch-Color Mappings & 1906–1960s & unclassified & discarded & none; no usable pitch-hue mappings \\
I.J. Belmont Equal-Division Spectral System & 1944 & F1 Spectral Analogy & retained & pitch-class-to-color; chromatic; 12 keys \\
August Aeppli Diatonic System & pre-1986 & F1 Spectral Analogy & retained & pitch-class-to-color; diatonic (8-note); 7 keys \\
Galeyev-Vanechkina / Prometei Institute Schema & 1975 & F2 Circle of Fifths & discarded & key-to-color; keys not normalizable to 12 pitch classes \\
Steve Zieverink Physical Wavelength Transposition & 2004 & F1 Spectral Analogy & retained & pitch-class-to-color; chromatic; 12 keys \\
Jack Ox Virtual Color Organ & 1977–present & F2 Circle of Fifths & discarded & key-to-color (with timbre extensions); keys not normalizable to 12 pitch classes \\
Katherine Lubar Itten Color Wheel Application & 2004 & F2 Circle of Fifths & discarded & color-interval analogy; not pitch-class mapping \\
Japanese Gogyo  Adopted Wuxing System & 6th century CE onward & F3 Cosmological & discarded & pitch-class-to-color (implicit via pentatonic); keys not normalizable to 12 pitch classes \\
Korean Ohaeng  Adopted Wuxing System & Traditional & F3 Cosmological & discarded & pitch-class-to-color (implicit via pentatonic); keys not normalizable to 12 pitch classes \\
Lawrence Marks Pitch-Brightness Correspondence & 1974/1975 & F7 Psychophysical & discarded & pitch-to-brightness/lightness only; not hue \\
Timothy Hubbard Pitch-Lightness Study & 1996 & F7 Psychophysical & discarded & pitch-to-brightness/lightness only; not hue \\
Mondloch and Maurer Developmental Pitch-Lightness & 2004 & F7 Psychophysical & discarded & pitch-to-brightness/lightness only; not hue \\
Palmer, Schloss, Xu, and Prado-León Emotion-Mediated Study & 2013 & F4 Emotion-Mediated & discarded & emotion/genre-mediated; not direct pitch-class \\
Palmer, Langlois, and Schloss Replication & 2016 & F4 Emotion-Mediated & discarded & emotion/genre-mediated; not direct pitch-class \\
Hamilton-Fletcher et al. Equiluminant Hue Study & 2017 & F7 Psychophysical & discarded & pitch-to-hue (controlled) \\
Whiteford et al. Genre-Dependent Music-Color & 2018 & F7 Psychophysical & discarded & emotion/genre-mediated; not direct pitch-class \\
Anikin and Johansson IAT Pitch-Hue Study & 2019 & F7 Psychophysical & discarded & pitch-to-hue (null result) \\
Spence and Di Stefano Comprehensive Review & 2022 & F7 Psychophysical & discarded & review (no novel mappings); no usable pitch-hue mappings \\
Loconsole et al. Tortoise Pitch-Lightness Study & 2025 & F7 Psychophysical & discarded & pitch-to-brightness/lightness only; not hue \\
Sean Day Compiled Synesthesia Databases & 2004–2016 & F5 Synesthetic & discarded & synesthesia aggregate/individual; not fixed system \\
Novich, Cheng, and Eagleman Synesthesia Battery & 2011 & F5 Synesthetic & discarded & synesthesia classification; no usable pitch-hue mappings \\
Gregersen et al. Absolute Pitch and Synesthesia & 2013 & F5 Synesthetic & discarded & synesthesia aggregate/individual; not fixed system \\
Friedrich Mahling Historical Bibliography & 1926 & unclassified & discarded & bibliography; no usable pitch-hue mappings \\
Neil Harbisson Sonochromatic Scale & 2003–present & F1 Spectral Analogy & discarded & continuous/microtone; not discrete pitch-class \\
Chroma-Notes / Boomwhackers Educational System & \textasciitilde{}1995 & F8 Pedagogical & retained & pitch-class-to-color; chromatic; 7 keys \\
Figurenotes Accessibility System & 1996 & F8 Pedagogical & retained & pitch-class-to-color; diatonic; 7 keys \\
EyeMusic Sensory Substitution Device & 2014 & F8 Pedagogical & discarded & string/timbre-to-color; not pitch-class \\
Chromatone Open-Source System & 2020s & F1 Spectral Analogy & retained & pitch-class-to-color; chromatic; 12 keys \\
Fred Collopy 'Three Centuries of Color Scales' Chart & 2009 & unclassified & discarded & comparative visualization; no usable pitch-hue mappings \\
Jörg Jewanski 'Ist C = Rot?' Doctoral Dissertation & 1999 & unclassified & discarded & historical compilation; no usable pitch-hue mappings \\
Kenneth Peacock 'Instruments to Perform Color-Music' & 1988 & unclassified & discarded & historical survey of color organs; no usable pitch-hue mappings \\
José Luis Caivano 'Color and Sound: Physical and Psychophysical Relations' & 1994 & unclassified & discarded & structural analysis; no usable pitch-hue mappings \\
Louis Bertrand Castel Ocular Harpsichord & 1725 & F1 Spectral Analogy & retained & pitch-class-to-color; chromatic; 12 keys; snowball addition \\
George Field Chromatic Equivalents & 1816 & F1 Spectral Analogy & retained & pitch-class-to-color; diatonic; 7 keys; snowball addition \\
Alexander Wallace Rimington Colour-Organ & 1895 & F1 Spectral Analogy & retained & pitch-class-to-color; chromatic; 12 keys; snowball addition \\
David A. Wells Proportional Frequency System & 1980 & F1 Spectral Analogy & retained & pitch-class-to-color; chromatic; 12 keys; snowball addition \\
Ralph W. Pridmore Wavelength-Frequency Transposition & 1992 & F1 Spectral Analogy & retained & pitch-class-to-color; chromatic; 12 keys; snowball addition \\
Soundmap (Mobile Application) & 2010s & F1 Spectral Analogy & retained & pitch-class-to-color; chromatic; 12 keys; snowball addition \\
\end{longtable}
\end{footnotesize}
\clearpage
\section*{S2. Family agreement uncertainty and the chance baseline}
\begin{table}[htbp]\centering
\caption{Table S2a. Family agreement scores with bootstrap 95\%
confidence intervals (10{,}000 draws, resamples constrained to at
least two distinct systems) and family-label permutation $p$-values
(10{,}000 permutations; $p$ is the proportion of label shuffles
producing a family score $\ge$ the observed one). A global
heterogeneity permutation test on the size-weighted variance of
family means is non-significant ($T_{\mathrm{obs}} = 0.0015$,
$p = .81$).}
\begin{tabular}{@{}lccccc@{}}
\toprule
\textbf{Family} & $N$ & \textbf{Agreement} & \textbf{95\% CI} & \textbf{Perm.\ $p$} & \textbf{Null mean (SD)} \\
\midrule
F8 Pedagogical & 3 & 0.348 & [0.301, 0.712] & 0.33 & 0.323 (0.147) \\
F5 Synesthetic & 5 & 0.276 & [0.196, 0.441] & 0.40 & 0.285 (0.100) \\
F1 Spectral Analogy & 24 & 0.267 & [0.229, 0.353] & 0.10 & 0.244 (0.017) \\
F6 Esoteric & 3 & 0.201 & [0.153, 0.260] & 0.80 & 0.324 (0.147) \\
F3 Cosmological & 4 & 0.187 & [0.178, 0.720] & 0.91 & 0.298 (0.118) \\
\bottomrule\end{tabular}\end{table}

\begin{table}[htbp]\centering
\caption{Table S2b. Simulated chance distribution of the agreement
score: ensembles of $n$ systems of 12 colors drawn uniformly from
the sRGB gamut, passed through the identical CIEDE2000 pipeline
(2{,}000 ensembles per $n$).}
\begin{tabular}{@{}ccc@{}}
\toprule
\textbf{Ensemble size $n$} & \textbf{Mean agreement} & \textbf{95\% interval} \\
\midrule
3 & 0.227 & [0.151, 0.408] \\
5 & 0.217 & [0.172, 0.310] \\
15 & 0.212 & [0.192, 0.240] \\
24 & 0.212 & [0.196, 0.231] \\
\bottomrule\end{tabular}\end{table}
\paragraph{Sensitivity check (31 retained systems).}
Recomputing the flat analyses on all 31 retained systems in the final
catalogue yields a mean per-pitch agreement of 0.237
(range 0.201--0.310),
Rayleigh $R = 0.60$, and a semitone--hue distance
correlation of $r = 0.49$
($p = 3.1e-05$), leaving all conclusions unchanged.
\clearpage
\section*{S3. Pedagogical consensus versus canonical spectral order}
To test whether the Pedagogical family's convergence simply reuses
rainbow order, the F8 consensus hues were compared against the
canonical spectral-order mapping (equally spaced hues in chromatic
pitch order). Over the seven pitch classes shared by all three F8
systems, the mean angular distance is 54.1$^\circ$;
allowing a free global rotation of the spectral template (hue origin
is arbitrary), the minimum mean distance is 35.1$^\circ$.
For calibration, the Spectral Analogy family (F1), whose systems
explicitly encode rainbow order, scores
39.9$^\circ$ unrotated and 35.8$^\circ$
after optimal rotation over all 12 pitch classes. The F8 consensus is
therefore about as close to spectral order (after rotation) as the
self-declared spectral family itself: the pedagogical convergence is
spectrally anchored, consistent with the reading in the main text
that educators converge on a spectrally ordered, affectively
comfortable mapping rather than on an arbitrary shared palette.
\clearpage
\section*{S4. Sentence-transformer probing grids}
\begin{table}[htbp]\centering
\caption{Table S4a. Full 3 (model) $\times$ 6 (descriptor level) probing grid, primary (LLM-generated, color-blocklisted) descriptor set. $p$-values from condition-specific label-shuffle permutation; $q$-values Benjamini--Hochberg corrected across the 18 cells. Bold: $p < .05$ before correction.}
\label{tab:s4a}
\begin{tabular}{@{}llccc@{}}
\toprule
\textbf{Model} & \textbf{Descriptor level} & \textbf{Geometry $\rho$} & \textbf{Perm.\ $p$} & \textbf{$q_{\mathrm{BH}}$} \\
\midrule
MiniLM & Bare names & -0.115 & 0.942 & 0.973 \\
MiniLM & Musical context & -0.039 & 0.590 & 0.817 \\
MiniLM & Perceptual & -0.040 & 0.631 & 0.817 \\
MiniLM & Timbral & -0.156 & 0.973 & 0.973 \\
MiniLM & Connotative & -0.049 & 0.645 & 0.817 \\
MiniLM & Solf\`ege & 0.019 & 0.344 & 0.774 \\
MPNet & Bare names & 0.055 & 0.226 & 0.582 \\
MPNet & Musical context & 0.156 & 0.060 & 0.251 \\
MPNet & Perceptual & \textbf{0.203} & \textbf{0.048} & 0.251 \\
MPNet & Timbral & -0.130 & 0.916 & 0.973 \\
MPNet & Connotative & 0.124 & 0.124 & 0.373 \\
MPNet & Solf\`ege & 0.000 & 0.441 & 0.794 \\
Nomic & Bare names & 0.173 & 0.070 & 0.251 \\
Nomic & Musical context & -0.032 & 0.516 & 0.817 \\
Nomic & Perceptual & \textbf{0.269} & \textbf{0.044} & 0.251 \\
Nomic & Timbral & \textbf{0.368} & \textbf{0.007} & 0.132 \\
Nomic & Connotative & -0.060 & 0.681 & 0.817 \\
Nomic & Solf\`ege & 0.005 & 0.402 & 0.794 \\
\bottomrule\end{tabular}\end{table}
\begin{table}[htbp]\centering
\caption{Table S4b. The same grid for the factual descriptor set (deterministically derived from music-theory constants). Bold: $p < .05$ before correction.}
\label{tab:s4b}
\begin{tabular}{@{}llccc@{}}
\toprule
\textbf{Model} & \textbf{Descriptor level} & \textbf{Geometry $\rho$} & \textbf{Perm.\ $p$} & \textbf{$q_{\mathrm{BH}}$} \\
\midrule
MiniLM & Bare names & -0.115 & 0.940 & 0.940 \\
MiniLM & Musical context & -0.036 & 0.573 & 0.737 \\
MiniLM & Perceptual & \textbf{0.164} & \textbf{0.033} & 0.155 \\
MiniLM & Timbral & \textbf{0.265} & \textbf{0.048} & 0.174 \\
MiniLM & Connotative & -0.128 & 0.850 & 0.900 \\
MiniLM & Solf\`ege & 0.049 & 0.249 & 0.499 \\
MPNet & Bare names & 0.055 & 0.221 & 0.497 \\
MPNet & Musical context & \textbf{0.334} & \textbf{0.012} & 0.139 \\
MPNet & Perceptual & -0.004 & 0.475 & 0.701 \\
MPNet & Timbral & \textbf{0.197} & \textbf{0.034} & 0.155 \\
MPNet & Connotative & -0.081 & 0.722 & 0.835 \\
MPNet & Solf\`ege & \textbf{0.230} & \textbf{0.015} & 0.139 \\
Nomic & Bare names & 0.173 & 0.058 & 0.175 \\
Nomic & Musical context & 0.109 & 0.129 & 0.331 \\
Nomic & Perceptual & -0.008 & 0.481 & 0.701 \\
Nomic & Timbral & -0.011 & 0.506 & 0.701 \\
Nomic & Connotative & -0.031 & 0.502 & 0.701 \\
Nomic & Solf\`ege & -0.089 & 0.742 & 0.835 \\
\bottomrule\end{tabular}\end{table}
\clearpage
\section*{S5. EPA Procrustes statistics}
\begin{table}[htbp]\centering
\caption{Table S5. EPA Procrustes statistics. Observed disparity $= 0.559$; permutation null (10{,}000 shuffles of the correspondence assignment): mean $= 0.833$, SD $= 0.106$; permutation $p = 0.027$. The sensorimotor configuration is the 3-component PCA (68.3\% of variance) of the 11-dimensional profiles of the 10 aligned dimensions. Dimension-wise Pearson correlations between matched coordinates after alignment are given below.}
\begin{tabular}{@{}lcc@{}}
\toprule
\textbf{EPA dimension} & \textbf{$r$} & \textbf{$p$} \\
\midrule
Evaluation & 0.74 & 0.014 \\
Potency    & 0.47 & 0.166 \\
Activity   & 0.37 & 0.293 \\
\bottomrule\end{tabular}\end{table}
\clearpage
\section*{S6. The coded CMC database}
\begin{footnotesize}
\begin{longtable}{@{}p{4.6cm}cp{2.6cm}p{1.9cm}p{4.4cm}@{}}
\caption{Table S6. The coded database of 34 documented cross-modal correspondence pairings: direction, literature-derived effect size (Cohen's $d$, with plausible range), mechanism proposed in the source literature, and key sources.}\label{tab:s6}\\
\toprule
\textbf{Pairing} & \textbf{Dir.} & \textbf{$d$ [range]} & \textbf{Mechanism} & \textbf{Key sources} \\
\midrule\endfirsthead
\toprule
\textbf{Pairing} & \textbf{Dir.} & \textbf{$d$ [range]} & \textbf{Mechanism} & \textbf{Key sources} \\
\midrule\endhead
\bottomrule\endfoot
pitch height $\to$ lightness & + & 1.00 [0.80, 1.20] & structural & Marks1974; Marks1987; Ward2006; Spence2011 \\
pitch height $\to$ spatial elevation & + & 0.80 [0.60, 1.00] & semantic & Spence2011; Rusconi2006 \\
pitch height $\to$ size & - & 0.65 [0.50, 0.80] & statistical & Mondloch2004; Spence2011 \\
pitch height $\to$ angularity & + & 0.50 [0.30, 0.70] & structural & Parise2012 \\
pitch height $\to$ hue & ? & 0.20 [0.05, 0.35] & unknown & Ward2006; Whiteford2018 \\
loudness $\to$ brightness & + & 0.40 [0.20, 0.60] & structural & Marks1987 \\
loudness $\to$ size & + & 0.60 [0.40, 0.80] & statistical & Smith1992 \\
tempo $\to$ saturation & + & 0.50 [0.30, 0.70] & emotion & Palmer2013 \\
tempo $\to$ lightness & + & 0.50 [0.30, 0.70] & emotion & Palmer2013; Collier2001 \\
mode major $\to$ hue warmth & + & 0.60 [0.40, 0.80] & emotion & Palmer2013 \\
spectral centroid $\to$ color temperature & + & 0.45 [0.25, 0.65] & structural & Giannakis2001; Palmer2013 \\
consonance $\to$ pleasantness & + & 0.55 [0.35, 0.75] & emotion & Lahdelma2016 \\
sweetness $\to$ roundness & + & 0.65 [0.45, 0.85] & statistical & Spence2011 \\
bitterness $\to$ angularity & + & 0.55 [0.35, 0.75] & statistical & Spence2011 \\
sourness $\to$ angularity & + & 0.40 [0.20, 0.60] & statistical & Velasco2016 \\
taste intensity $\to$ saturation & + & 0.50 [0.30, 0.70] & structural & Spence2011 \\
saltiness $\to$ whiteness & + & 0.35 [0.20, 0.50] & semantic & Wan2014 \\
umami $\to$ brownness & + & 0.30 [0.15, 0.45] & semantic & Spence2011 \\
smoothness $\to$ lightness & + & 0.40 [0.20, 0.60] & structural & Ludwig2013 \\
hardness $\to$ loudness & + & 0.40 [0.20, 0.60] & statistical & Walker2012 \\
roughness $\to$ spectral roughness & + & 0.50 [0.30, 0.70] & structural & Guest2002 \\
temperature warm $\to$ hue warmth & + & 0.45 [0.25, 0.65] & semantic & Deroy2013 \\
heaviness $\to$ pitch height & - & 0.35 [0.15, 0.55] & statistical & Walker2012 \\
odor pleasantness $\to$ lightness & + & 0.40 [0.25, 0.55] & hedonic & Kemp1997 \\
odor intensity $\to$ saturation & + & 0.38 [0.18, 0.58] & hedonic & Gilbert2008 \\
citrus $\to$ yellow & + & 0.42 [0.22, 0.62] & semantic & Schifferstein1999 \\
mint $\to$ green & + & 0.36 [0.20, 0.52] & semantic & Schifferstein1999 \\
pitch direction $\to$ brightness & + & 0.35 [0.15, 0.55] & structural & Collier2001 \\
rhythmic complexity $\to$ angularity & + & 0.30 [0.15, 0.45] & emotion & Palmer2013 \\
timbre brightness $\to$ hue warmth & + & 0.44 [0.24, 0.64] & structural & Palmer2013 \\
brightness $\to$ loudness & + & 0.28 [0.12, 0.44] & structural & Marks1987 \\
saturation $\to$ taste intensity & + & 0.46 [0.26, 0.66] & statistical & Spence2011 \\
redness $\to$ sweetness & + & 0.41 [0.21, 0.61] & semantic & Spence2015 \\
greenness $\to$ sourness & + & 0.37 [0.17, 0.57] & semantic & Velasco2016 \\
\end{longtable}
\end{footnotesize}
\clearpage
\section*{S7. Consolidated multiplicity table}
\begin{table}[htbp]\centering
\caption{Table S7. Consolidated multiplicity table: all primary and
battery-level tests with Benjamini--Hochberg FDR-corrected
$q$-values.}
\begin{tabular}{@{}p{9.5cm}cc@{}}
\toprule
\textbf{Test} & \textbf{$p$} & \textbf{$q_{\mathrm{BH}}$} \\
\midrule
Semitone-hue correlation H2 (permutation) & 0.0005 & 0.005 \\
Interval-variance correlation H\_b6 & 0.0014 & 0.007 \\
Timbral embedding alignment (Nomic, permutation) & 0.0073 & 0.024 \\
Rayleigh non-uniformity (Monte Carlo) & 0.0103 & 0.026 \\
EPA Procrustes alignment (permutation) & 0.0268 & 0.045 \\
Effect size vs sensorimotor distance & 0.0269 & 0.045 \\
Agreement-consistency H\_b9 & 0.0325 & 0.046 \\
Family heterogeneity: F1 label-permutation & 0.1011 & 0.126 \\
Family heterogeneity: F8 label-permutation & 0.3330 & 0.370 \\
Family heterogeneity: F5 label-permutation & 0.3973 & 0.397 \\
\bottomrule\end{tabular}\end{table}
\clearpage
\section*{S8. Mean hue distance by musical interval}
\begin{table}[htbp]\centering
\caption{Table S8. Mean hue distance (degrees on the colour wheel)
by musical interval, pooled across all 25 retained systems. Circular
pitch-class distance makes the perfect fourth and perfect fifth
equivalent (both 5 semitones); their rows are identical by
construction.}
\begin{tabular}{@{}lccc@{}}
\toprule
\textbf{Interval} & \textbf{Semitones} &
  \textbf{Mean hue dist.\ ($^\circ$)} & \textbf{SD ($^\circ$)} \\
\midrule
Minor second   & 1 & 36.9 & 44.6 \\
Major second   & 2 & 53.3 & 44.3 \\
Minor third    & 3 & 59.1 & 38.3 \\
Major third    & 4 & 70.5 & 34.9 \\
Perfect fourth & 5 & 80.7 & 26.1 \\
Tritone        & 6 & 85.7 & 22.0 \\
Perfect fifth  & 7 & 80.7 & 26.1 \\
\bottomrule
\end{tabular}
\end{table}

\clearpage
\section*{S9. Corpus-level hypothesis battery (full 10 tests)}
\begin{table}[htbp]\centering
\caption{Table S9. Corpus-level hypothesis test results. The Account
column attributes each hypothesis to the theoretical account it
operationalises; only the Statistical rows bear directly on the
environmental co-occurrence account.}
\begin{tabular}{@{}p{0.9cm}p{4.6cm}p{1.7cm}p{1.9cm}p{3.2cm}@{}}
\toprule
\textbf{Hyp.} & \textbf{Description} & \textbf{Account} &
  \textbf{Result} & \textbf{Key statistic} \\
\midrule
$H_{b1}$ & Lightness increases with chromatic pitch order
  & Structural & Fail & $r = -.53$, $p = .076$ \\
$H_{b2}$ & Agreement is higher for high-chroma assignments
  & Hedonic & Fail & $r = .55$, $p = .063$ \\
$H_{b3}$ & Consonant intervals show larger hue distance than dissonant
  & Structural & Fail & $t = 1.05$, $p = .34$ \\
$H_{b4}$ & Hue distribution non-uniform
  & Statistical & \textbf{Pass} & Rayleigh $R = 0.66$, $p = .003$ \\
$H_{b5}$ & Sharps differ from naturals in lightness or hue
  & Cultural & Fail & $p = .53$ / $.95$ \\
$H_{b6}$ & Interval size predicts between-system variance
  & Statistical & \textbf{Pass} & $r = -.94$, $p = .001$ \\
$H_{b7}$ & Perfect fifth $<$ tritone hue distance
  & Statistical & Inconclusive & $80.7^\circ < 85.7^\circ$ (untested) \\
$H_{b8}$ & Primary triad (C--E--G) hues $\sim$120$^\circ$ apart
  & Structural & Fail & mean $66.4^\circ$, dev.\ $53.6^\circ$ \\
$H_{b9}$ & Agreement correlates with consistency
  & Statistical & \textbf{Pass} & $r = -.62$, $p = .03$ \\
$H_{b10}$ & Tradition-specific fifths differ from pooled
  & Cultural & Fail & $t = 0.15$, $p = .89$ \\
\bottomrule
\end{tabular}
\end{table}

\clearpage
\section*{S10. Per-family harmonic structure}
\begin{figure}[htbp]\centering
\includegraphics[width=0.85\textwidth]{P1_03v2_family_ranking.png}
\caption{Figure S10. Per-family harmonic structure analysis for three
families (Esoteric F6, Spectral Analogy F1, Synesthetic F5). Left:
harmonic structure score (mean hue distance for dissonant minus
consonant intervals); all three families assign more similar hues to
consonant than dissonant intervals. Right: Spearman $\rho$ between
semitone distance and mean hue distance; Spectral Analogy and
Synesthetic families show substantially stronger interval-to-hue
structure ($\rho \approx 0.6$) than the Esoteric family.}
\end{figure}

\clearpage
\section*{S11. Structural theory pairwise alignment heatmap}
\begin{figure}[htbp]\centering
\includegraphics[width=0.75\textwidth]{P1_05_theory_heatmap.png}
\caption{Figure S11. Pairwise alignment heatmap for the six
structural theories. Cell colour encodes mean angular distance between
predicted hue assignments. IRP and FRI form the tightest pair
(40.1$^\circ$), confirming that they capture the same structural
intuition through different mathematical formalisms.}
\end{figure}

\clearpage
\section*{S12. Structural theory predicted hue assignments}
\begin{figure}[htbp]\centering
\includegraphics[width=0.75\textwidth]{P1_05_theory_wheel.png}
\caption{Figure S12. Predicted hue assignments for the six structural
theories overlaid on the global consensus wheel. IRP is the closest
match to the observed historical consensus, though it diverges
substantially in the middle register (D--G$\sharp$).}
\end{figure}

\clearpage
\section*{S13. Per-model language-model predicted hue maps}
\begin{figure}[htbp]\centering
\includegraphics[width=\textwidth]{P1_07_lm_wheels.png}
\caption{Figure S13. Predicted hue assignments from the three
language models (MiniLM, MPNet, Nomic) under LLM-generated timbral
descriptors. Nomic shows the most structured prediction (warm reds at
lower pitches, cooler hues at higher pitches), consistent with its
higher geometry alignment ($\rho = .368$, $p = .007$).}
\end{figure}

\clearpage
\section*{S14. Cross-linguistic and emotion-mediation control panels}
\begin{figure}[htbp]\centering
\includegraphics[width=\textwidth]{P1_07v2_crossling_emotion.png}
\caption{Figure S14. Left: geometry alignment across five naming
conventions (English, Japanese, German, solf\`{e}ge with context, bare
solf\`{e}ge); pitch--hue geometry is weak across all naming conventions,
arguing against name-specific artefacts. Right: the small positive
Spearman correlation ($\rho = .139$, $p < .001$) between pitch--colour
embedding proximity and shared emotional valence is consistent with
the hedonic account.}
\end{figure}

\clearpage
\section*{S15. Complete per-system mapping panel by reasoning family}
\begin{figure}[htbp]\centering
\includegraphics[width=\textwidth]{P1_families_panel.png}
\caption{Figure S15. Within-family mapping panel. Each row shows one
historical system's colour assignments for all 12 pitch classes (C to
B), grouped by reasoning family (F1--F8). Spectral Analogy (F1) is
dominated by rainbow order; Synesthetic (F5) tends toward warm hues
at low pitches and cool hues at high; Pedagogical (F8) shows the
tightest within-family clustering.}
\end{figure}

\end{document}